\documentclass[11pt]{article}

\usepackage[a4paper,margin=1.6cm]{geometry}
\usepackage{titlesec}
\usepackage{enumitem}
\usepackage[hidelinks]{hyperref}
\usepackage{graphicx}
\usepackage{array}
\usepackage{xeCJK}

\linespread{0.95} % 行距略缩
\setlength{\parindent}{0pt}
\setlength{\parskip}{0pt}
\setlist[itemize]{leftmargin=*, topsep=2pt, itemsep=0pt, parsep=0pt, partopsep=0pt}
\titleformat{\section}{\large\bfseries}{}{0em}{}[\titlerule]
\titlespacing*{\section}{0pt}{0.3em}{0.15em}

\begin{document}

%---------- 基本信息 + 照片（垂直居中对齐） ----------
\noindent
\begin{minipage}[c]{0.76\textwidth}
\textbf{\large 李星宇} \\
性别：男 \quad 出生年月：2002.06 \\
电话：18361227161 \\
邮箱：1690855076@qq.com \\
籍贯：江苏省徐州市 \\
GitHub：\href{https://github.com/xqsrdtd}{github.com/xqsrdtd}
\end{minipage}
\hfill
\begin{minipage}[c]{0.19\textwidth}
\centering
% \includegraphics[width=2.35cm,height=3.1cm]{photo_1.jpg}
% \includegraphics[width=2.35cm,height=3.1cm]{photo_3.jpg}
\IfFileExists{photo_4.png}{%
  \includegraphics[width=2.4cm,height=3cm]{photo_4.png}%
}{%
  \fbox{\parbox[c][3cm][c]{2.4cm}{\centering Photo}}%
}
\end{minipage}

\vspace{-2pt}

%----------教育背景----------
\section{教育背景}

\textbf{哈尔滨工业大学（深圳）}｜机械工程硕士（机器人与先进制造学院） \hfill 2024.09 – 至今 \\
研究方向：仿生机器人运动控制、强化学习

\textbf{北京航空航天大学}｜机器人工程 本科 \hfill 2020.09 – 2024.07

\vspace{-2pt}

%----------求职意向----------
\section{求职意向}

强化学习算法实习生 / 具身智能算法实习生（机器人控制方向）

\vspace{1pt}

%----------项目经历----------
\section{核心项目经历}

\textbf{人形机器人延迟鲁棒性强化学习（Isaac Gym, PPO）}｜独立项目 \hfill 2025.09 – 至今
\begin{itemize}
\item 构建 Baseline 与 Domain Randomization 对照实验，通过引入 DOF lag 实现执行器时延鲁棒策略学习
\item 训练阶段随机化 0–40 timesteps 时延，测试阶段推广至 OOD 场景（最高 80 timesteps）
\item 实验结果：lag=80 时鲁棒策略成功率 90\%（Baseline 为 0\%）；lag=0 时达 99.3\%
\item 完成多随机种子验证、实验自动化 Pipeline、可视化分析与 Isaac Gym→MuJoCo sim2sim 迁移
\item 项目代码：\href{https://github.com/xqsrdtd/agibot_x1_dof_lag}{github.com/xqsrdtd/agibot_x1_dof_lag}
\end{itemize}

\textbf{基于语言的机械臂规划与控制（VoxPoser 改进）}｜项目负责人 \hfill 2024.01 – 2024.07
\begin{itemize}
\item 基于 VoxPoser 3D 体素规划框架轻量化适配与算法改进，实现自然语言指令驱动机械臂末端轨迹规划
\item 融合 A* 提出体素衰减累积成本机制，引入动态权重策略 + Bézier 曲线平滑，提升路径质量与执行稳定性
\item 在 RLBench/CoppeliaSim 完成 4 类任务测试（10 次重复），整体成功率 85\%
\item 完成算法部署与 Aubo E5 实体机械臂实验验证
\end{itemize}

\textbf{两栖六足仿生机器人（CPG 控制）}｜项目负责人 \hfill 2024.09 – 至今
\begin{itemize}
\item 设计陆水两栖结构，实现单套推进器多模态仿生运动
\item 基于 CPG 振荡器网络实现步态生成与相位协调控制
\item 通过结构-控制联合调参提升运动稳定性
\end{itemize}

\textbf{全国大学生机器人大赛 RR 机器人研发}｜技术负责人 \hfill 2022.09 – 2023.09
\begin{itemize}
\item 主导机器人总体方案设计与两轮迭代，负责核心机构设计及机构/硬件/控制联合调试
\item 组织团队完成原型开发、赛场测试与策略优化，沉淀从需求拆解到系统落地的完整工程流程
\item 获全国大学生机器人大赛国家级二等奖
\end{itemize}

\vspace{-2pt}

%----------论文发表----------
\section{论文发表}

{\small \textbf{Cockroach’s Turning Strategy Enhanced Hexapod Robot with Flexible Torso}\hspace{-0.05em}（IROS 2025，二作）}
\begin{itemize}
\item 负责仿生步态控制算法设计，基于蟑螂转向行为提取步态相位关系与占空比参数
\item 构建六足机器人转向控制策略（gait-based 与 mix-based），实现步态与躯干弯曲协同控制
\item 实验验证：相比纯步态策略，混合策略转弯半径降低约 27.4\%，转向速度提升约 40\%
\end{itemize}
\vspace{-2pt}

%----------实习经历----------
\section{实习经历}

\textbf{遨博（北京）智能科技有限公司}｜技术实习生 \hfill 2023.07 – 2023.08
\begin{itemize}
\item 参与机械臂控制系统开发、装配调试与性能测试
\item 熟悉工业机器人软硬件联调流程
\end{itemize}

\vspace{-2pt}

%----------技能----------
\section{技能}

\textbf{强化学习与仿真：}PPO（主用），SAC，TD3；Domain Randomization；Isaac Gym，MuJoCo，RLBench；sim2sim / sim2real

\textbf{编程与工程：}Python（PyTorch），C，MATLAB；实验框架搭建、数据分析与可视化

\textbf{机器人系统：}机器人学、机器人控制技术、执行器延迟建模、系统集成与调参优化

\vspace{-2pt}

%----------荣誉----------
\section{竞赛与荣誉}

全国大学生机器人大赛 国家级二等奖 \\
全国部分地区大学生物理竞赛 国家级一等奖 \\
校级工程设计表达竞赛一等奖、学科竞赛奖学金一等奖等

\end{document}
